\documentclass[11pt]{article}
\usepackage[margin=1in]{geometry}
\usepackage{amsmath,amssymb,amsthm,mathrsfs}
\usepackage{graphicx}
\usepackage[hidelinks]{hyperref}

\newtheorem{theorem}{Theorem}
\newtheorem{remark}{Remark}

\newcommand{\R}{\mathbb{R}}
\newcommand{\N}{\mathbb{N}}
\newcommand{\Z}{\mathbb{Z}}
\newcommand{\E}{\mathbb{E}}
\newcommand{\cK}{\mathscr{K}}
\newcommand{\norm}[1]{\left\lVert #1 \right\rVert}
\newcommand{\abs}[1]{\left\lvert #1 \right\rvert}
\newcommand{\Lp}{L^p}
\newcommand{\Hs}{H^s_p}
\newcommand{\UMD}{\mathrm{UMD}}
\newcommand{\Rad}{\mathrm{Rad}}
\newcommand{\lesssimC}{\lesssim}

\title{Unsigned Difference Norms in the Banach Function Space Case}
\author{Agent lane 14}
\date{June 28, 2026}

\begin{document}
\maketitle

\section*{Source and Target}

The source paper is Nick Lindemulder, \emph{Difference Norms for
Vector-Valued Bessel Potential Spaces with an Application to Pointwise
Multipliers}, arXiv:1505.08039.  Corollary 4.9 proves a square-function
difference norm characterization for \(H^s_p(\R^d,w;X)\) when \(X\) is a
\(\UMD\) Banach function space.  Immediately after it, Remark 4.10 asks whether
one may replace the signed averaged differences \(K_m(2^{-j},f)\) in the
right hand side by unsigned difference averages \(d^m_K(2^{-j},f)\).

\begin{center}
\includegraphics[width=0.96\textwidth]{figures/open_problem_crop.png}
\end{center}

There is an evident typo in the displayed definition in Remark 4.10: the scale
\(t\) is missing from the difference.  The surrounding sentence
\(\abs{K_m(t,f)}\le d^m_K(t,f)\) only makes sense for
\[
  D^m_K(t,f)(x):=\int_{\R^d} K(h)\,\abs{\Delta^m_{t h}f(x)}\,dh.
\]
The theorem below answers the question for this intended scaled definition,
under the same standing hypotheses as Corollary 4.9.

\section*{Proof Intuition}

The signed square function in Corollary 4.9 is controlled by the source paper
through a Littlewood-Paley decomposition \(f=\sum S_n f\).  The unsigned
quantity is sublinear rather than linear, but in a Banach function space the
lattice modulus and the \(\ell^2\)-square function let us keep the same
decomposition.  High-frequency pieces are controlled by the source
\(\ell^2\)-boundedness of positive convolution kernels.  Low-frequency pieces
are controlled exactly as in the proof of Theorem 4.1(ii): for each fixed
\(h\), the difference \(\Delta^m_{2^{-j}h}\) is represented by the multiplier
family of Lemma 4.8, and the integrability assumption on \(K\) absorbs the
factor \((1+\abs h)^{(d+3)m}\).  The source inequality
\(\abs{K_m(t,f)}\le D^m_K(t,f)\) gives the converse side.

\begin{theorem}
Let \(X\) be a \(\UMD\) Banach function space, let \(s>0\),
\(p\in(1,\infty)\), \(w\in A_p(\R^d)\), and let \(m\in\N\) with \(m>s\).
Assume that
\[
  K\in \cK(\R^d)^+\cap L^1(\R^d,(1+\abs{\cdot})^{(d+3)m})
\]
and that \(K^{\Delta^m}\) satisfies the Tauberian condition of
Lindemulder's Theorem 4.1(i), equivalently the condition used in Corollary
4.9.  Then, for every \(f\in L^p(\R^d,w;X)\),
\[
\norm{f}_{H^s_p(\R^d,w;X)}
\simeq
\norm{f}_{L^p(\R^d,w;X)}
+
\norm{
  \left(\sum_{j=1}^{\infty}
    \abs{2^{js}D^m_K(2^{-j},f)}^2
  \right)^{1/2}
}_{L^p(\R^d,w;X)}.
\]
The constants depend only on the parameters in the source corollary and on
\(\int K(h)(1+\abs h)^{(d+3)m}\,dh\).
\end{theorem}

\begin{proof}
Write
\[
  U_J(f):=
  \norm{
    \left(\sum_{j=1}^{J}
      \abs{2^{js}D^m_K(2^{-j},f)}^2
    \right)^{1/2}
  }_{L^p(\R^d,w;X)}.
\]
It is enough to prove \(U_J(f)\lesssimC \norm{f}_{H^s_p(\R^d,w;X)}\)
uniformly in \(J\), since the opposite estimate follows from Corollary 4.9 of
the source paper and the lattice inequality
\[
  \abs{K_m(t,f)}=
  \abs{\int_{\R^d}K(h)\Delta^m_{t h}f\,dh}
  \le
  \int_{\R^d}K(h)\abs{\Delta^m_{t h}f}\,dh
  =D^m_K(t,f).
\]

Fix a Littlewood-Paley decomposition \((S_n)_{n\ge 0}\) as in the source
paper and put \(f_n=S_nf\) for \(n\ge0\), \(f_n=0\) for \(n<0\).  For each
\(j\ge1\), \(f=\sum_{n\in\Z}f_{n+j}\) in the relevant distribution sense.
Since \(D^m_K(t,\cdot)\) is subadditive and positive, the lattice Minkowski
inequality gives
\[
  U_J(f)
  \le
  \sum_{n\in\Z}
  \norm{
    \left(\sum_{j=1}^{J}
      \abs{2^{js}D^m_K(2^{-j},f_{n+j})}^2
    \right)^{1/2}
  }_{L^p(\R^d,w;X)}.
\]

First let \(n>0\).  From
\[
  \abs{\Delta^m_{t h}g}
  \le
  \sum_{\ell=0}^{m}{m\choose \ell}\abs{g(\,\cdot +(m-\ell)t h)}
\]
we get
\[
  D^m_K(t,g)
  \le
  \sum_{\ell=0}^{m} {m\choose \ell} P_{\ell,t}\abs g,
\]
where \(P_{\ell,t}\) is either multiplication by \(\int K\) or convolution
with a positive dilation/reflection of \(K\).  These kernels form a bounded
set in \(\cK(\R^d)\).  Proposition 4.11 of the source paper says that such
bounded sets of convolution kernels are \(\ell^2\)-bounded on
\(L^p(\R^d,w;X)\) when \(X\) is a \(\UMD\) Banach function space.  Hence
\[
\begin{aligned}
&\norm{
    \left(\sum_{j=1}^{J}
      \abs{2^{js}D^m_K(2^{-j},f_{n+j})}^2
    \right)^{1/2}
  }_{L^p(w;X)}
\\
&\qquad\lesssim
 2^{-ns}
 \norm{
    \left(\sum_{j=1}^{J}
      \abs{2^{(n+j)s}f_{n+j}}^2
    \right)^{1/2}
  }_{L^p(w;X)}
 \lesssim
 2^{-ns}\norm{f}_{H^s_p(\R^d,w;X)}.
\end{aligned}
\]
The last step is the Khintchine-Maurey theorem together with the randomized
Littlewood-Paley equivalence (13) in the source paper.

Now let \(n\le0\).  By lattice Minkowski and Khintchine-Maurey,
\[
\begin{aligned}
&\norm{
    \left(\sum_{j=1}^{J}
      \abs{2^{js}D^m_K(2^{-j},f_{n+j})}^2
    \right)^{1/2}
  }_{L^p(w;X)}
\\
&\qquad\le
 \int_{\R^d} K(h)
 \norm{
    \left(\sum_{j=1}^{J}
      \abs{2^{js}\Delta^m_{2^{-j}h}f_{n+j}}^2
    \right)^{1/2}
  }_{L^p(w;X)}\,dh
\\
&\qquad\simeq
 \int_{\R^d} K(h)
 \norm{\sum_{j=1}^{J}\varepsilon_j
      2^{js}\Delta^m_{2^{-j}h}f_{n+j}}
      _{L^2(\Omega;L^p(w;X))}\,dh .
\end{aligned}
\]
The proof of Theorem 4.1(ii) in the source paper applies Lemma 4.8 to the
last Rademacher norm and yields, for every fixed \(h\),
\[
 \norm{\sum_{j=1}^{J}\varepsilon_j
      2^{js}\Delta^m_{2^{-j}h}f_{n+j}}
      _{L^2(\Omega;L^p(w;X))}
 \lesssim
 2^{n(m-s)}(1+\abs h)^{(d+3)m}
 \norm{f}_{H^s_p(\R^d,w;X)}.
\]
Integrating in \(h\) is allowed by the assumed moment condition on \(K\), so
\[
  \norm{
    \left(\sum_{j=1}^{J}
      \abs{2^{js}D^m_K(2^{-j},f_{n+j})}^2
    \right)^{1/2}
  }_{L^p(w;X)}
  \lesssim
  2^{n(m-s)}\norm{f}_{H^s_p(\R^d,w;X)}.
\]

Summing the estimates over \(n>0\) and \(n\le0\) gives two convergent
geometric series, since \(s>0\) and \(m-s>0\).  Thus
\[
  U_J(f)\lesssim \norm{f}_{H^s_p(\R^d,w;X)}
\]
uniformly in \(J\).  Letting \(J\to\infty\) proves the desired upper estimate
by monotone convergence in the Banach lattice \(L^p(\R^d,w;X)\).  Together
with the already proved lower estimate this gives the asserted norm
equivalence.
\end{proof}

\section*{Verification and Novelty Notes}

No computation is used.  The proof reduces the open estimate to three source
ingredients: Corollary 4.9, Proposition 4.11, and the fixed-step multiplier
estimate in Lemma 4.8 as used in the proof of Theorem 4.1(ii).

Bounded duplicate and novelty checks were performed in the local run indexes
for arXiv id 1505.08039, the phrase
\(\texttt{R-bdd\_averaging\_operators}\), averaging operators, and
unsigned/difference-modulus keywords.  No existing packet for this question
was found.  A bounded web search on June 28, 2026 for exact phrases from
Remark 4.10 and for \(\texttt{d\_K}\), \(\texttt{K\_m}\), and the paper title
did not identify a later answer.  Novelty confidence is therefore
\(\texttt{bounded-search likely new}\), not a complete bibliographic audit.

\section*{Limitations and Review Recommendation}

This packet answers the intended scaled version of Remark 4.10.  Read
literally, the source display of \(d^m_K(t,f)\) omits \(t\) from
\(\Delta^m_{t h}\); that literal formula cannot be the intended one because
it would not imply the displayed domination \(\abs{K_m(t,f)}\le d^m_K(t,f)\).

Human review should focus on the high-frequency use of Proposition 4.11 for
the positive convolution majorants and on the transfer of the source
low-frequency Rademacher estimate through Khintchine-Maurey and lattice
Minkowski.

\begin{thebibliography}{9}
\bibitem{Lindemulder}
N. Lindemulder,
\emph{Difference Norms for Vector-Valued Bessel Potential Spaces with an
Application to Pointwise Multipliers},
arXiv:1505.08039.
\end{thebibliography}

\end{document}
